\chapter{Chapter06 — Gowers}
%%% generated by the blueprint pipeline from FABL.Chapter06.Gowers
\section{Overview}

Declarations of the Lean module \texttt{FABL.Chapter06.Gowers}.

\section{Declarations}

\begin{definition}[Real-valued functions on pairs of binary cubes]
\lean{FABL.GowersFamily}
\statementsource{Analysis-of-Boolean-Functions-by-Ryan-ODonnell}{pdf-ffbe9c7b3306-p168-b004}
\label{FABL.GowersFamily}
\leanok
\uses{FABL.F₂Cube}
For natural numbers $n$ and $k$, a Gowers family is a real-valued function of a pair
$(s,x)$, where $s\in\mathbb{F}_2^k$ and $x\in\mathbb{F}_2^n$.
\end{definition}

\begin{definition}[Gowers cube point]
\lean{FABL.gowersCubePoint}
\statementsource{Analysis-of-Boolean-Functions-by-Ryan-ODonnell}{pdf-ffbe9c7b3306-p168-b004}
\label{FABL.gowersCubePoint}
\leanok
\uses{FABL.F₂Cube}
Let $x\in\mathbb{F}_2^n$, let $y=(y_i)_{i=1}^k$ be a $k$-tuple of points in
$\mathbb{F}_2^n$, and let $s\in\mathbb{F}_2^k$. The associated cube point is
$x+\sum_{i=1}^k s_i y_i$, with addition and scalar multiplication in $\mathbb{F}_2^n$.
\end{definition}

\begin{definition}[Gowers inner average]
\lean{FABL.gowersInner}
\statementsource{Analysis-of-Boolean-Functions-by-Ryan-ODonnell}{pdf-ffbe9c7b3306-p168-b004}
\label{FABL.gowersInner}
\leanok
\uses{FABL.F₂Cube, FABL.GowersFamily, FABL.gowersCubePoint}
For natural numbers $n$ and $k$, let $F$ assign a real number $F(s,z)$ to each pair
$s\in\mathbb{F}_2^k$ and $z\in\mathbb{F}_2^n$. Define the Gowers inner average of $F$ to
be the average over all functions $y\colon\mathbb{F}_2^k\to\mathbb{F}_2^n$ and all
$x\in\mathbb{F}_2^n$ of the product $\prod_{s\in\mathbb{F}_2^k}
F\bigl(s,\,g(x,y,s)\bigr)$, where $g(x,y,s)$ is the point specified by the Gowers cube
construction.
\end{definition}

\begin{definition}[Diagonal Gowers inner average]
\lean{FABL.gowersDiagonalInner}
\statementsource{Analysis-of-Boolean-Functions-by-Ryan-ODonnell}{pdf-ffbe9c7b3306-p168-b004}
\label{FABL.gowersDiagonalInner}
\leanok
\uses{FABL.F₂Cube, FABL.gowersInner}
For natural numbers $n$ and $k$ and a real-valued function
$f\colon\mathbb{F}_2^n\to\mathbb{R}$, define its diagonal Gowers inner average of order
$k$ to be the Gowers inner average of the family that takes the value $f(x)$ at every
pair $(s,x)$.
\end{definition}

\begin{definition}[Gowers norm]
\lean{FABL.gowersNorm}
\statementsource{Analysis-of-Boolean-Functions-by-Ryan-ODonnell}{pdf-ffbe9c7b3306-p168-b004}
\label{FABL.gowersNorm}
\leanok
\uses{FABL.F₂Cube, FABL.gowersDiagonalInner}
For natural numbers $n$ and $k$ and a function $f\colon\mathbb{F}_2^n\to\mathbb{R}$,
define the Gowers norm of order $k$ by raising the diagonal Gowers inner product of $f$
of order $k$ to the power $1/2^k$.
\end{definition}

\begin{definition}[Constant binary vector of length one]
\lean{FABL.gowersBit}
\label{FABL.gowersBit}
\leanok
\uses{FABL.F₂Cube, FABL.𝔽₂}
For $b\in\mathbb{F}_2$, define the length-one binary vector whose sole coordinate is
$b$.
\end{definition}

\begin{definition}[Binary pair as a length-two vector]
\lean{FABL.gowersPair}
\label{FABL.gowersPair}
\leanok
\uses{FABL.F₂Cube, FABL.𝔽₂}
For $a,b\in\mathbb{Z}/2\mathbb{Z}$, let $v\colon\{0,1\}\to\mathbb{Z}/2\mathbb{Z}$ be the
vector with $v(0)=a$ and $v(1)=b$.
\end{definition}

\begin{theorem}[Insertion identity for Gowers cube points]
\lean{FABL.gowersCubePoint_insertNth}
\label{FABL.gowersCubePoint_insertNth}
\leanok
\uses{FABL.F₂Cube, FABL.gowersCubePoint, FABL.𝔽₂}
\difficulty{3}
Let $n,k\in\mathbb N$, let $i\in\{0,\ldots,k\}$, and take $x,h\in\mathbb F_2^n$, vectors
$y_0,\ldots,y_{k-1}\in\mathbb F_2^n$, a bit $b\in\mathbb F_2$, and
$s=(s_0,\ldots,s_{k-1})\in\mathbb F_2^k$. If $h$ is inserted into the list $y$ and $b$
into $s$ at position $i$, then the resulting cube point equals the cube point formed
from base point $x+bh$ and the original $y,s$; that is, both equal $x+\sum_{j<k}s_jy_j+b
h$.
\end{theorem}

\begin{theorem}[Product over a binary cube split]
\lean{FABL.prod_f₂Cube_succ}
\leanok
\uses{FABL.F₂Cube, FABL.𝔽₂}
\difficulty{4}
Let $M$ be a commutative monoid, let $k$ be a natural number, and let $f$ assign an
element of $M$ to each $(k+1)$-tuple over the two-element field. The product of $f$ over
all such tuples is the product of its values on the tuples whose last coordinate is $0$
and its values on the tuples whose last coordinate is $1$.
\end{theorem}

\begin{theorem}[Product split over a set and its complement]
\lean{FABL.prod_ite_mem_compl}
\label{FABL.prod_ite_mem_compl}
\leanok
\difficulty{4}
Let $\iota$ be a finite type, let $M$ be a commutative monoid, let $A$ be a subset of
$\iota$, and let $f,g:\iota\to M$. The product over all $i\in\iota$ that uses $f(i)$ on
$A$ and $g(i)$ outside $A$ equals the product of $f$ over $A$ times the product of $g$
over the complement of $A$.
\end{theorem}

\begin{theorem}[Expectation over a product of finite types]
\lean{FABL.expect_prod_type}
\label{FABL.expect_prod_type}
\leanok
\difficulty{3}
Let $alpha$ and $beta$ be finite types and let $f:alphatimesbetatobbr$. The
uniform expectation of $f$ over pairs $(a,b)$ equals the iterated uniform expectation,
first over $binbeta$ and then over $ainalpha$.
\end{theorem}

\begin{theorem}[Expectation of a translated product]
\lean{FABL.expect_mul_add_eq_mul_expect}
\label{FABL.expect_mul_add_eq_mul_expect}
\leanok
\uses{FABL.F₂Cube}
\difficulty{4}
Let $n$ be a nonnegative integer, and let $a,b$ be real-valued functions on the binary
cube $\mathbb F_2^n$. If $x$ and $h$ are independent and uniformly distributed on that
cube, then $\mathbb E[a(x)b(x+h)]=(\mathbb E a)(\mathbb E b)$.
\end{theorem}

\begin{definition}[Gowers half product]
\lean{FABL.gowersHalfProductAt}
\label{FABL.gowersHalfProductAt}
\leanok
\uses{FABL.F₂Cube, FABL.GowersFamily, FABL.gowersCubePoint, FABL.𝔽₂}
Let $F$ be a real-valued family on pairs consisting of a binary vector of length $k+1$
and a point of the $n$-dimensional binary cube. Choose a coordinate
$i\in\{0,\ldots,k\}$, a bit $b\in\mathbb{F}_2$, a point $x$ of the $n$-dimensional
binary cube, and a $k$-tuple $y$ of such points. The Gowers half product is the product,
over all binary vectors $s$ of length $k$, of $F$ evaluated at the vector obtained by
inserting $b$ into $s$ at coordinate $i$ and at the Gowers cube point associated to
$x,y,s$.
\end{definition}

\begin{definition}[Gowers half expectation]
\lean{FABL.gowersHalfExpectationAt}
\label{FABL.gowersHalfExpectationAt}
\leanok
\uses{FABL.F₂Cube, FABL.GowersFamily, FABL.gowersHalfProductAt, FABL.𝔽₂}
Let $F$ be a real-valued family on pairs consisting of a binary vector of length $k+1$
and a point of the $n$-dimensional binary cube. Choose a coordinate
$i\in\{0,\ldots,k\}$, a bit $b\in\mathbb{F}_2$, and a $k$-tuple $y$ of points of the
$n$-dimensional binary cube. The Gowers half expectation is the uniform average, over
all points $x$ of that cube, of the Gowers half product at $x$ for the chosen $F,i,b,y$.
\end{definition}

\begin{definition}[Half expectation at the last coordinate]
\lean{FABL.gowersHalfExpectation}
\label{FABL.gowersHalfExpectation}
\leanok
\uses{FABL.F₂Cube, FABL.GowersFamily, FABL.gowersHalfExpectationAt, FABL.𝔽₂}
For a family $F:([?]$ and a bit $b\in\mathbb F_2$, with $y$ a $k$-tuple of points of
$\mathbb F_2^n$, define its half expectation to be the half-expectation functional for
$F$ at the last coordinate of the $(k+1)$-cube, evaluated at $b$ and $y$.
\end{definition}

\begin{theorem}[Factorization of a Gowers cube product]
\lean{FABL.gowersProduct_insertNth}
\label{FABL.gowersProduct_insertNth}
\leanok
\uses{FABL.F₂Cube, FABL.GowersFamily, FABL.gowersCubePoint, FABL.gowersCubePoint_insertNth, FABL.gowersHalfProductAt, FABL.𝔽₂}
\difficulty{4}
Let $F$ assign a real-valued function on $\mathbb{F}_2^n$ to each point of
$\mathbb{F}_2^{k+1}$, let $i\in\{0,\ldots,k\}$, and let $x,h\in\mathbb{F}_2^n$ and
$y_1,\ldots,y_k\in\mathbb{F}_2^n$. The product over $s\in\mathbb{F}_2^{k+1}$ of $F(s)$
evaluated at $x+\sum_{j=1}^{k+1}s_j y'_j$, where $y'$ is obtained by inserting $h$ into
$y$ at position $i$, equals the product of the two half-cube products: the slice with
inserted coordinate $0$ based at $x$, times the slice with inserted coordinate $1$ based
at $x+h$.
\end{theorem}

\begin{theorem}[Splitting the inner product at a coordinate]
\lean{FABL.gowersInner_split_at}
\label{FABL.gowersInner_split_at}
\leanok
\uses{FABL.F₂Cube, FABL.GowersFamily, FABL.expect_mul_add_eq_mul_expect, FABL.expect_prod_type, FABL.gowersCubePoint, FABL.gowersHalfExpectationAt, FABL.gowersHalfProductAt, FABL.gowersInner, FABL.gowersProduct_insertNth}
\difficulty{5}
Let $n,k\in\mathbb N$, let $F$ be a real-valued family indexed by $F_2^{k+1}$ with
values that are functions on $F_2^n$, and choose a coordinate $i\in[k+1]$. Its Gowers
inner product equals the average over $y\in(F_2^n)^k$ of the product of two slice
expectations. For each bit $b\in\mathbb F_2$, the corresponding slice expectation is the
average over $x\in F_2^n$ of the product over $s\in F_2^k$ of $F$ evaluated at the point
formed by inserting $b$ into coordinate $i$ of $s$, and at $x+\sum_j s_jy_j$.
\end{theorem}

\begin{theorem}[Successor-dimensional inner product decomposition]
\lean{FABL.gowersInner_succ}
\label{FABL.gowersInner_succ}
\leanok
\uses{FABL.F₂Cube, FABL.GowersFamily, FABL.gowersHalfExpectation, FABL.gowersInner, FABL.gowersInner_split_at}
\difficulty{1}
Let $n,k\in\mathbb N$ and let $F$ be a real-valued family indexed by $F_2^{k+1}$ with
values that are functions on $F_2^n$. Its Gowers inner product equals the average over
$y\in(F_2^n)^k$ of the product of the two expectations obtained by fixing the last cube
coordinate to $0$ and to $1$. Each expectation averages over $x\in F_2^n$ the product of
the family values over the remaining $k$ cube coordinates, evaluated at the
corresponding points $x+\sum_j s_jy_j$.
\end{theorem}

\begin{definition}[Coordinate duplication of a family]
\lean{FABL.gowersDuplicateAt}
\label{FABL.gowersDuplicateAt}
\leanok
\uses{FABL.GowersFamily, FABL.𝔽₂}
Let $F$ assign a real-valued function on $\mathbb F_2^n$ to each point of the
$k$-dimensional cube, choose a coordinate $i$ and a bit $b\in\mathbb F_2$. Define a new
family by replacing coordinate $i$ of each cube point by $b$ before evaluating $F$.
\end{definition}

\begin{definition}[Duplication at the last coordinate]
\lean{FABL.gowersLastDuplicate}
\label{FABL.gowersLastDuplicate}
\leanok
\uses{FABL.GowersFamily, FABL.gowersDuplicateAt, FABL.𝔽₂}
Let $F$ assign a real-valued function on $\mathbb F_2^n$ to each point of the
$(k+1)$-dimensional cube, and let $b\in\mathbb F_2$. Define a new family by replacing
the last coordinate of each cube point by $b$ before evaluating $F$.
\end{definition}

\begin{theorem}[Last-coordinate duplication identity]
\lean{FABL.gowersLastDuplicate_apply}
\label{FABL.gowersLastDuplicate_apply}
\leanok
\uses{FABL.F₂Cube, FABL.GowersFamily, FABL.gowersDuplicateAt, FABL.gowersLastDuplicate, FABL.𝔽₂}
\difficulty{1}
Let $n,k\in\mathbb N$, let $F$ be a real-valued family indexed by $F_2^{k+1}$ with
values that are functions on $F_2^n$, let $b\in\mathbb F_2$, and let $s\in F_2^{k+1}$.
If the last-coordinate duplication of $F$ at $b$ is evaluated at $s$, the result is $F$
evaluated at the point whose first $k$ coordinates are those of $s$ and whose last
coordinate is $b$.
\end{theorem}

\begin{theorem}[Duplicate-coordinate inner product formula]
\lean{FABL.gowersInner_duplicateAt}
\label{FABL.gowersInner_duplicateAt}
\leanok
\uses{FABL.F₂Cube, FABL.GowersFamily, FABL.gowersDuplicateAt, FABL.gowersHalfExpectationAt, FABL.gowersHalfProductAt, FABL.gowersInner, FABL.gowersInner_split_at, FABL.𝔽₂}
\difficulty{4}
Let $F$ be a Gowers family of order $k+1$, choose a coordinate $i\in\mathrm{Fin}(k+1)$,
and choose $b\in\mathbb F_2$. The Gowers inner product of the family obtained by
replacing coordinate $i$ by $b$ equals the average over $y:\mathrm{Fin}(k)\to\mathbb
F_2^n$ of the square of the corresponding half expectation:
$\operatorname{Inner}(\operatorname{duplicate}_{i,b}F)=\mathbb
E_y\bigl[\operatorname{HalfExpectation}_{i,b}(F;y)^2\bigr]$.
\end{theorem}

\begin{theorem}[Nonnegativity of a duplicated Gowers inner product]
\lean{FABL.gowersInner_duplicateAt_nonneg}
\label{FABL.gowersInner_duplicateAt_nonneg}
\leanok
\uses{FABL.GowersFamily, FABL.gowersDuplicateAt, FABL.gowersInner, FABL.gowersInner_duplicateAt, FABL.𝔽₂}
\difficulty{1}
Let $n,k\in\mathbb N$, let $F$ assign a real-valued function on the $n$-dimensional
Boolean cube to each point of the $k$-dimensional Boolean cube, and choose a coordinate
$i$ and a bit $b\in\mathbb F_2$. Replace the input to $F$ by the point obtained by
setting coordinate $i$ to $b$. The resulting Gowers inner product, formed by averaging
over $y\in(F_2^n)^{k+1}$ and $x\in F_2^n$ the product over all $s\in F_2^{k+1}$ of the
family value at $x+\sum_j s_jy_j$, is nonnegative.
\end{theorem}

\begin{theorem}[Nonnegativity of the diagonal inner product]
\lean{FABL.gowersDiagonalInner_nonneg}
\label{FABL.gowersDiagonalInner_nonneg}
\leanok
\uses{FABL.F₂Cube, FABL.GowersFamily, FABL.gowersDiagonalInner, FABL.gowersDuplicateAt, FABL.gowersInner, FABL.gowersInner_duplicateAt_nonneg}
\difficulty{4}
For every nonnegative integer $k$ and every real-valued function $f$ on the
$n$-dimensional Boolean cube, the diagonal Gowers inner product of order $k+1$ is
nonnegative: $0\leq\operatorname{diag}_{k+1}(f)$.
\end{theorem}

\begin{theorem}[Nonnegativity of the Gowers norm]
\lean{FABL.gowersNorm_nonneg}
\label{FABL.gowersNorm_nonneg}
\leanok
\uses{FABL.F₂Cube, FABL.gowersDiagonalInner_nonneg, FABL.gowersNorm}
\difficulty{4}
Let $n$ be a nonnegative integer, let $k$ be a positive integer, and let $f\colon\mathbb
F_2^n\to\mathbb R$. Define $\|f\|_{U^k}$ as the $2^k$-th root of $\mathbb
E_{x,y_1,\ldots,y_k}\prod_{s\in\mathbb F_2^k}f(x+\sum_{i=1}^k s_i y_i)$, with the
average uniform over $x,y_1,\ldots,y_k\in\mathbb F_2^n$. Then $0\leq\|f\|_{U^k}$.
\end{theorem}

\begin{theorem}[Power identity for the Gowers norm]
\lean{FABL.gowersNorm_pow}
\label{FABL.gowersNorm_pow}
\leanok
\uses{FABL.F₂Cube, FABL.gowersDiagonalInner, FABL.gowersDiagonalInner_nonneg, FABL.gowersNorm}
\difficulty{3}
Let $n$ be a nonnegative integer, let $k$ be a positive integer, and let $f\colon\mathbb
F_2^n\to\mathbb R$. Define $\|f\|_{U^k}$ as the $2^k$-th root of $\mathbb
E_{x,y_1,\ldots,y_k}\prod_{s\in\mathbb F_2^k}f(x+\sum_{i=1}^k s_i y_i)$, with the
average uniform over $x,y_1,\ldots,y_k\in\mathbb F_2^n$. Then $\|f\|_{U^k}^{2^k}=\mathbb
E_{x,y_1,\ldots,y_k}\prod_{s\in\mathbb F_2^k}f(x+\sum_{i=1}^k s_i y_i)$.
\end{theorem}

\begin{theorem}[Coordinatewise Gowers Cauchy–Schwarz inequality]
\lean{FABL.gowers_cauchy_schwarz_at}
\label{FABL.gowers_cauchy_schwarz_at}
\leanok
\uses{FABL.F₂Cube, FABL.GowersFamily, FABL.gowersDuplicateAt, FABL.gowersHalfExpectationAt, FABL.gowersInner, FABL.gowersInner_duplicateAt, FABL.gowersInner_split_at}
\difficulty{5}
Let $F$ assign a real-valued function on $\mathbb{F}_2^n$ to each $s\in\mathbb{F}_2^k$,
and fix a coordinate $i\in\{0,\ldots,k-1\}$. The Gowers inner product of $F$ is at most
the product of the square roots of the Gowers inner products of the two families
obtained by replacing the $i$th coordinate of every index by $0$ and by $1$,
respectively.
\end{theorem}

\begin{theorem}[Last-coordinate Gowers Cauchy–Schwarz inequality]
\lean{FABL.gowers_cauchy_schwarz_last}
\label{FABL.gowers_cauchy_schwarz_last}
\leanok
\uses{FABL.GowersFamily, FABL.gowersInner, FABL.gowersLastDuplicate, FABL.gowers_cauchy_schwarz_at}
\difficulty{1}
Let $F$ assign a real-valued function on $\mathbb{F}_2^n$ to each
$s\in\mathbb{F}_2^{k+1}$. Its Gowers inner product is at most the product of the square
roots of the Gowers inner products of the two families obtained by replacing the last
coordinate of every index by $0$ and by $1$, respectively.
\end{theorem}

\begin{definition}[Lifting a family along an embedding]
\lean{FABL.gowersLiftFamilyAlong}
\label{FABL.gowersLiftFamilyAlong}
\leanok
\uses{FABL.GowersFamily}
Let $e$ be an embedding of the $r$ coordinate positions into the $k$ coordinate
positions, and let $F$ assign a real-valued function on $\mathbb F_2^n$ to each point of
the $r$-dimensional cube. Define a family on the $k$-dimensional cube by restricting
each input to the coordinates selected by $e$ and then evaluating $F$.
\end{definition}

\begin{theorem}[Duplicating a coordinate after lifting]
\lean{FABL.gowersDuplicateAt_liftFamilyAlong}
\label{FABL.gowersDuplicateAt_liftFamilyAlong}
\leanok
\uses{FABL.F₂Cube, FABL.GowersFamily, FABL.gowersDuplicateAt, FABL.gowersLiftFamilyAlong, FABL.𝔽₂}
\difficulty{4}
Let $e:\mathrm{Fin}(r+1)\hookrightarrow\mathrm{Fin}(k)$ be an embedding, let $F$ be a
Gowers family of order $r+1$, and let $b\in\mathbb F_2$. Duplicating the lifted family
at the coordinate $e(r)$ with value $b$ equals the lift along the embedding obtained by
restricting $e$ to its first $r$ coordinates, applied to the family that appends $b$ to
each input: $\operatorname{duplicate}_{e(r),b}(\operatorname{lift}_e
F)=\operatorname{lift}_{e|_{\mathrm{Fin}(r)}}(s\mapsto F(s\mathbin{\|}b))$.
\end{theorem}

\begin{definition}[Recursive Cauchy–Schwarz bound]
\lean{FABL.gowersCSBound}
\label{FABL.gowersCSBound}
\leanok
\uses{FABL.GowersFamily, FABL.gowersDiagonalInner}
For a family $F$ indexed by the $r$-dimensional cube with values in real-valued
functions on $\mathbb F_2^n$, define a real quantity recursively. At $r=0$, it is the
diagonal Gowers inner average of the function assigned to the unique point of the
zero-dimensional cube. At $r+1$, it is the product of the square roots of the quantities
at level $r$ obtained by appending, respectively, the bit $0$ and the bit $1$ to each
$r$-dimensional input.
\end{definition}

\begin{theorem}[Lifted inner product bounded by the Cauchy–Schwarz bound]
\lean{FABL.gowersInner_liftFamilyAlong_le_csBound}
\label{FABL.gowersInner_liftFamilyAlong_le_csBound}
\leanok
\uses{FABL.GowersFamily, FABL.gowersCSBound, FABL.gowersDiagonalInner, FABL.gowersDuplicateAt_liftFamilyAlong, FABL.gowersInner, FABL.gowersLiftFamilyAlong, FABL.gowers_cauchy_schwarz_at}
\difficulty{5}
Let $n,k,r\in\mathbb N$, let $e:[r]\hookrightarrow[k]$ be an injection, and let $F$ be a
real-valued family indexed by the $r$-dimensional Boolean cube with values that are
functions on the $n$-dimensional Boolean cube. Lift $F$ to dimension $k$ by evaluating
it on the restriction of each $k$-cube point along $e$. The Gowers inner product of this
lifted family is at most the recursively defined Cauchy–Schwarz bound: at dimension zero
this is the diagonal inner product, and at each successor dimension it is the product of
the square roots of the bounds for the two slices obtained by fixing the last coordinate
to $0$ and to $1$.
\end{theorem}

\begin{theorem}[Square root of a product of diagonal averages]
\lean{FABL.sqrt_prod_gowersDiagonal_rpow}
\label{FABL.sqrt_prod_gowersDiagonal_rpow}
\leanok
\uses{FABL.F₂Cube, FABL.GowersFamily, FABL.gowersDiagonalInner, FABL.gowersDiagonalInner_nonneg}
\difficulty{4}
Let $k>0$, let $r$ be a nonnegative integer, and let $F$ assign a real-valued function
on the $n$-dimensional Boolean cube to each point of the $r$-dimensional Boolean cube.
For each such point $s$, let $D_k(F_s)$ be the average, over a uniformly chosen base
point and $k$ directions in the $n$-dimensional Boolean cube, of the product of $F_s$
evaluated at the vertices of the resulting $k$-dimensional cube. Then
\[\sqrt{\prod_{s\in\{0,1\}^r} D_k(F_s)^{1/2^r}}=\prod_{s\in\{0,1\}^r}D_k(F_s)^{1/2^{r+1}}.\]
\end{theorem}

\begin{theorem}[Product formula for the Cauchy–Schwarz bound]
\lean{FABL.gowersCSBound_eq_prod_rpow}
\label{FABL.gowersCSBound_eq_prod_rpow}
\leanok
\uses{FABL.F₂Cube, FABL.GowersFamily, FABL.gowersCSBound, FABL.gowersDiagonalInner, FABL.prod_f₂Cube_succ, FABL.sqrt_prod_gowersDiagonal_rpow}
\difficulty{4}
Let $k$ be a positive integer, let $r$ be a nonnegative integer, and let $F$ assign a
real-valued function on the $n$-dimensional Boolean cube to each point of the
$r$-dimensional Boolean cube. The recursively defined Gowers Cauchy–Schwarz bound
satisfies $\displaystyle
\operatorname{CS}_k(F)=\prod_{s\in\{0,1\}^r}\operatorname{diag}_k(F(s))^{1/2^r}$, where
$\operatorname{diag}_k(f)$ denotes the diagonal Gowers inner product of $f$.
\end{theorem}

\begin{theorem}[Cauchy–Schwarz bound as a product of norms]
\lean{FABL.gowersCSBound_eq_prod_norm}
\label{FABL.gowersCSBound_eq_prod_norm}
\leanok
\uses{FABL.GowersFamily, FABL.gowersCSBound, FABL.gowersCSBound_eq_prod_rpow, FABL.gowersNorm}
\difficulty{1}
Let $n,k$ be natural numbers with $k>0$, and let $F$ assign a real-valued function on
$\mathbb F_2^n$ to each $s\in\mathbb F_2^k$. The recursively defined Cauchy–Schwarz
bound for $F$ at level $k$ equals $\prod_{s\in\mathbb F_2^k}\|F(s)\|_{U^k}$, where
$\|f\|_{U^k}$ is the $2^k$-th root of the diagonal Gowers inner product of order $k$ of
$f$.
\end{theorem}

\begin{theorem}[Gowers Cauchy–Schwarz inequality]
\lean{FABL.gowers_cauchy_schwarz}
\label{FABL.gowers_cauchy_schwarz}
\leanok
\uses{FABL.GowersFamily, FABL.gowersCSBound_eq_prod_norm, FABL.gowersInner, FABL.gowersInner_liftFamilyAlong_le_csBound, FABL.gowersLiftFamilyAlong, FABL.gowersNorm}
\difficulty{3}
Let $k$ be a positive integer and let $F$ assign a real-valued function on
$\mathbb{F}_2^n$ to each $s\in\mathbb{F}_2^k$. Its Gowers inner product is at most the
product, over all $s\in\mathbb{F}_2^k$, of the order-$k$ Gowers norms of the functions
$F(s)$.
\end{theorem}

\begin{theorem}[One-dimensional Gowers inner product]
\lean{FABL.gowersInner_one_eq}
\proofsource{Analysis-of-Boolean-Functions-by-Ryan-ODonnell}{pdf-ffbe9c7b3306-p168-b005}
\label{FABL.gowersInner_one_eq}
\leanok
\uses{FABL.F₂Cube, FABL.GowersFamily, FABL.gowersBit, FABL.gowersCubePoint, FABL.gowersHalfExpectation, FABL.gowersHalfExpectationAt, FABL.gowersHalfProductAt, FABL.gowersInner, FABL.gowersInner_succ, FABL.mean, FABL.𝔽₂}
\difficulty{4}
Let $n\in\mathbb N$ and let $F$ assign a real-valued function on $F_2^n$ to each point
of $F_2^1$. The Gowers inner product of $F$ equals the product of the uniform means of
its two slices, indexed by the one-coordinate points $0$ and $1$.
\end{theorem}

\begin{theorem}[Square of the first Gowers norm]
\lean{FABL.gowersNorm_one_sq}
\proofsource{Analysis-of-Boolean-Functions-by-Ryan-ODonnell}{pdf-ffbe9c7b3306-p168-b005}
\label{FABL.gowersNorm_one_sq}
\leanok
\uses{FABL.F₂Cube, FABL.gowersDiagonalInner, FABL.gowersInner_one_eq, FABL.gowersNorm, FABL.gowersNorm_pow, FABL.mean}
\difficulty{3}
Let $n$ be a nonnegative integer and let $f\colon\mathbb F_2^n\to\mathbb R$. Write
$\mathbb E_x f(x)$ for the uniform mean of $f$ on $\mathbb F_2^n$. Then the square of
the first Gowers norm of $f$ equals the square of its mean: $\|f\|_{U^1}^2=(\mathbb E_x
f(x))^2$.
\end{theorem}

\begin{theorem}[Half expectation as a convolution]
\lean{FABL.gowersHalfExpectation_two_eq_convolution}
\label{FABL.gowersHalfExpectation_two_eq_convolution}
\leanok
\uses{FABL.F₂Cube, FABL.GowersFamily, FABL.convolution, FABL.convolution_apply_add, FABL.gowersCubePoint, FABL.gowersHalfExpectation, FABL.gowersHalfExpectationAt, FABL.gowersHalfProductAt, FABL.gowersPair, FABL.prod_f₂Cube_succ, FABL.𝔽₂}
\difficulty{4}
Let $F$ be a family indexed by pairs of bits and taking real-valued functions on the
$n$-dimensional Boolean cube, let $b\in\mathbb F_2$, and let
$y:\mathrm{Fin}(1)\to\mathbb F_2^n$. The half expectation obtained by fixing the last
index to $b$ equals the normalized additive convolution of the functions indexed by
$(0,b)$ and $(1,b)$, evaluated at $y(0)$.
\end{theorem}

\begin{theorem}[Two-dimensional inner product and Fourier coefficients]
\lean{FABL.gowersInner_two_eq_fourier}
\label{FABL.gowersInner_two_eq_fourier}
\leanok
\uses{FABL.F₂Cube, FABL.GowersFamily, FABL.convolution, FABL.gowersHalfExpectation_two_eq_convolution, FABL.gowersInner, FABL.gowersInner_succ, FABL.gowersPair, FABL.vectorFourierCoeff, FABL.vectorFourierCoeff_convolution, FABL.vector_plancherel}
\difficulty{4}
Let $n\in\mathbb N$ and let $F$ be a real-valued family indexed by $F_2^2$ with values
that are functions on $F_2^n$. Its Gowers inner product equals the sum, over all $[?]$
\end{theorem}

\begin{theorem}[Fourth power of the order-two Gowers norm]
\lean{FABL.gowersNorm_two_pow_four_eq}
\proofsource{Analysis-of-Boolean-Functions-by-Ryan-ODonnell}{pdf-ffbe9c7b3306-p168-b006}
\label{FABL.gowersNorm_two_pow_four_eq}
\leanok
\uses{FABL.F₂Cube, FABL.gowersDiagonalInner, FABL.gowersInner_two_eq_fourier, FABL.gowersNorm, FABL.gowersNorm_pow, FABL.vectorFourierFourthMoment}
\difficulty{3}
For a function $f:\mathbb{F}_2^n\to\mathbb{R}$, the fourth power of its order-two Gowers
norm equals the sum of the fourth powers of all its Fourier coefficients:
$\|f\|_{U^2}^4=\sum_{\gamma\in\mathbb{F}_2^n}\widehat f(\gamma)^4$.
\end{theorem}

\begin{theorem}[Gowers norm of the constant function]
\lean{FABL.gowersNorm_one}
\label{FABL.gowersNorm_one}
\leanok
\uses{FABL.F₂Cube, FABL.gowersDiagonalInner, FABL.gowersInner, FABL.gowersNorm, FABL.gowersNorm_nonneg, FABL.gowersNorm_pow}
\difficulty{3}
Let $n$ be a nonnegative integer and let $k$ be a positive integer. The Gowers norm of
order $k$ of the constant function $1$ on $\mathbb F_2^n$ is $1$.
\end{theorem}

\begin{definition}[Monotonicity family]
\lean{FABL.gowersMonotonicityFamily}
\label{FABL.gowersMonotonicityFamily}
\leanok
\uses{FABL.F₂Cube, FABL.GowersFamily}
Let $f:\mathbb F_2^n\to\mathbb R$. Define a family indexed by the $(k+1)$-dimensional
cube that takes the value $f$ when the last coordinate is $0$ and the constant function
with value $1$ when the last coordinate is nonzero.
\end{definition}

\begin{theorem}[Inner product of the monotonicity family]
\lean{FABL.gowersInner_monotonicityFamily}
\label{FABL.gowersInner_monotonicityFamily}
\leanok
\uses{FABL.F₂Cube, FABL.gowersDiagonalInner, FABL.gowersHalfExpectation, FABL.gowersHalfExpectationAt, FABL.gowersHalfProductAt, FABL.gowersInner, FABL.gowersInner_succ, FABL.gowersMonotonicityFamily}
\difficulty{3}
Let $n,k\in\mathbb N$ and let $f:F_2^n\to\mathbb R$. Form a family indexed by $s\in
F_2^{k+1}$ that equals $f$ when the last coordinate of $s$ is $0$, and equals the
constant function $1$ when that coordinate is nonzero. Its Gowers inner product equals
the diagonal inner product in dimension $k$, namely the inner product of the family that
assigns $f$ to every point of $F_2^k$.
\end{theorem}

\begin{theorem}[Product of norms in a monotonicity family]
\lean{FABL.prod_gowersNorm_monotonicityFamily}
\label{FABL.prod_gowersNorm_monotonicityFamily}
\leanok
\uses{FABL.F₂Cube, FABL.card_f₂Cube, FABL.gowersMonotonicityFamily, FABL.gowersNorm, FABL.gowersNorm_one, FABL.prod_f₂Cube_succ}
\difficulty{4}
Let $n,k$ be natural numbers and let $f$ be a real-valued function on the
$n$-dimensional cube over the two-element field. Form the family on the
$(k+1)$-dimensional cube by taking the function $f$ when the last coordinate is $0$ and
the constant function $1$ when it is $1$. The product of the order-$(k+1)$ Gowers norms
of these functions over all cube indices equals the order-$(k+1)$ Gowers norm of $f$
raised to $2^k$.
\end{theorem}

\begin{theorem}[Monotonicity of the Gowers norm]
\lean{FABL.gowersNorm_mono}
\label{FABL.gowersNorm_mono}
\leanok
\uses{FABL.F₂Cube, FABL.gowersInner_monotonicityFamily, FABL.gowersMonotonicityFamily, FABL.gowersNorm, FABL.gowersNorm_nonneg, FABL.gowersNorm_pow, FABL.gowers_cauchy_schwarz, FABL.prod_gowersNorm_monotonicityFamily}
\difficulty{3}
Let $n$ be a nonnegative integer, let $k$ be a positive integer, and let $f\colon\mathbb
F_2^n\to\mathbb R$. Define $\|f\|_{U^j}$ as the $2^j$-th root of $\mathbb
E_{x,y_1,\ldots,y_j}\prod_{s\in\mathbb F_2^j}f(x+\sum_{i=1}^j s_i y_i)$, with the
average uniform over $x,y_1,\ldots,y_j\in\mathbb F_2^n$. Then
$\|f\|_{U^k}\leq\|f\|_{U^{k+1}}$.
\end{theorem}

\begin{definition}[Choice family on a finite set]
\lean{FABL.gowersChoiceFamily}
\label{FABL.gowersChoiceFamily}
\leanok
\uses{FABL.F₂Cube, FABL.GowersFamily}
Let $f,g:\mathbb F_2^n\to\mathbb R$ and let $A$ be a finite subset of the
$k$-dimensional cube. Define a family that takes the value $f$ at inputs in $A$ and the
value $g$ at inputs outside $A$.
\end{definition}

\begin{theorem}[Expansion of the diagonal inner product]
\lean{FABL.gowersDiagonalInner_add_expansion}
\label{FABL.gowersDiagonalInner_add_expansion}
\leanok
\uses{FABL.F₂Cube, FABL.gowersChoiceFamily, FABL.gowersCubePoint, FABL.gowersDiagonalInner, FABL.gowersInner, FABL.prod_ite_mem_compl}
\difficulty{4}
Let $f,g$ be real-valued functions on the $n$-dimensional Boolean cube, and let $k$ be a
nonnegative integer. The diagonal Gowers inner product of $f+g$ equals the sum, over
every subset $A$ of the $k$-dimensional Boolean cube, of the Gowers inner product of the
family that uses $f$ at cube indices in $A$ and $g$ at all other indices.
\end{theorem}

\begin{theorem}[Choice-family identity for Gowers norms]
\lean{FABL.sum_prod_gowersNorm_choice}
\label{FABL.sum_prod_gowersNorm_choice}
\leanok
\uses{FABL.F₂Cube, FABL.card_f₂Cube, FABL.gowersChoiceFamily, FABL.gowersNorm, FABL.prod_ite_mem_compl}
\difficulty{3}
Let $n,k$ be natural numbers and let $f,g:\mathbb F_2^n\to\mathbb R$, where $\mathbb
F_2$ is the two-element field. For each subset $A$ of $\mathbb F_2^k$, choose the
function equal to $f$ on indices in $A$ and to $g$ on the other indices; the sum over
all such $A$ of the product, over $s\in\mathbb F_2^k$, of the Gowers $k$-norm of the
chosen function at index $s$ equals $(\|f\|_{U^k}+\|g\|_{U^k})^{2^k}$.
\end{theorem}

\begin{theorem}[Triangle inequality for the Gowers norm]
\lean{FABL.gowersNorm_add_le}
\label{FABL.gowersNorm_add_le}
\leanok
\uses{FABL.F₂Cube, FABL.gowersChoiceFamily, FABL.gowersDiagonalInner, FABL.gowersDiagonalInner_add_expansion, FABL.gowersInner, FABL.gowersNorm, FABL.gowersNorm_nonneg, FABL.gowersNorm_pow, FABL.gowers_cauchy_schwarz, FABL.sum_prod_gowersNorm_choice}
\difficulty{4}
Let $n$ be a nonnegative integer and $k$ a positive integer, and let $f,g\colon\mathbb
F_2^n\to\mathbb R$. For $h\colon\mathbb F_2^n\to\mathbb R$, write $\|h\|_{U^k}$ for the
$2^k$-th root of $\mathbb E_{x,y_1,\ldots,y_k}\prod_{s\in\mathbb F_2^k}h(x+\sum_{i=1}^k
s_i y_i)$, where the average is uniform over all choices in $\mathbb F_2^n$. Then
$\|f+g\|_{U^k}\leq \|f\|_{U^k}+\|g\|_{U^k}$.
\end{theorem}

\begin{theorem}[Even powers and absolute values]
\lean{FABL.pow_two_pow_eq_abs_pow_two_pow}
\label{FABL.pow_two_pow_eq_abs_pow_two_pow}
\leanok
\difficulty{3}
For every positive integer $k$ and every real number $a$, one has $a^{2^k}=|a|^{2^k}$.
\end{theorem}

\begin{theorem}[Homogeneity of the diagonal inner product]
\lean{FABL.gowersDiagonalInner_smul}
\label{FABL.gowersDiagonalInner_smul}
\leanok
\uses{FABL.F₂Cube, FABL.card_f₂Cube, FABL.gowersCubePoint, FABL.gowersDiagonalInner, FABL.gowersInner}
\difficulty{4}
For every real number $a$ and every real-valued function $f$ on the $n$-dimensional
Boolean cube, the diagonal Gowers inner product of order $k$ satisfies
$\operatorname{diag}_k(af)=a^{2^k}\operatorname{diag}_k(f)$.
\end{theorem}

\begin{theorem}[Homogeneity of the Gowers norm]
\lean{FABL.gowersNorm_smul}
\label{FABL.gowersNorm_smul}
\leanok
\uses{FABL.F₂Cube, FABL.gowersDiagonalInner_smul, FABL.gowersNorm, FABL.gowersNorm_nonneg, FABL.gowersNorm_pow, FABL.pow_two_pow_eq_abs_pow_two_pow}
\difficulty{3}
Let $n$ be a nonnegative integer, let $k$ be a positive integer, let $a\in\mathbb R$,
and let $f\colon\mathbb F_2^n\to\mathbb R$. Then $\|af\|_{U^k}=|a|\,\|f\|_{U^k}$.
\end{theorem}

\begin{definition}[Gowers seminorm]
\lean{FABL.gowersSeminorm}
\label{FABL.gowersSeminorm}
\leanok
\uses{FABL.F₂Cube, FABL.gowersNorm, FABL.gowersNorm_add_le, FABL.gowersNorm_smul}
For a positive integer $k$, define the seminorm on real-valued functions on $\mathbb
F_2^n$ obtained from the $k$th Gowers norm.
\end{definition}

\begin{theorem}[Comparison with higher-order Gowers norms]
\lean{FABL.gowersNorm_two_le}
\label{FABL.gowersNorm_two_le}
\leanok
\uses{FABL.F₂Cube, FABL.gowersNorm, FABL.gowersNorm_mono}
\difficulty{4}
Let $n$ be a nonnegative integer, let $k$ be an integer with $2\leq k$, and let
$f\colon\mathbb F_2^n\to\mathbb R$. Then $\|f\|_{U^2}\leq\|f\|_{U^k}$.
\end{theorem}

\begin{theorem}[Vanishing Gowers norm forces vanishing]
\lean{FABL.eq_zero_of_gowersNorm_eq_zero}
\proofsource{Analysis-of-Boolean-Functions-by-Ryan-ODonnell}{pdf-ffbe9c7b3306-p169-b007}
\label{FABL.eq_zero_of_gowersNorm_eq_zero}
\leanok
\uses{FABL.F₂Cube, FABL.gowersNorm, FABL.gowersNorm_nonneg, FABL.gowersNorm_two_le, FABL.gowersNorm_two_pow_four_eq, FABL.vectorFourierCoeff, FABL.vectorFourierFourthMoment, FABL.vector_fourier_expansion}
\difficulty{5}
Let $n$ be a natural number, let $k\ge 2$, and let $f:(\mathbb F_2)^n\to\mathbb R$. If
the Gowers norm of order $k$ of $f$ is zero, then $f$ is identically zero.
\end{theorem}
